\section{All rational characters, the trace pairing, and finite windows}
\label{sec:trace-windows}

This section uses the following source abbreviations. C99 is Connes,
\emph{Trace formula in noncommutative geometry and the zeros of the
Riemann zeta function}, the supplied \texttt{main.tex}
(arXiv:math/9811068). CCM07 is Connes--Consani--Marcolli,
\emph{The Weil proof and the geometry of the adeles class space},
the supplied \texttt{Yuri70v2.tex} (arXiv:math/0703392).
CC20 is Connes--Consani, \emph{Weil positivity and Trace formula:
the archimedean place}, the supplied \texttt{weil-compo.tex}
(arXiv:2006.13771). The statement numbers below have been checked
against the theorem-counter declarations and the complete sequence
of numbered environments in these supplied sources. Source labels
are also supplied, to make the locators independent of typesetting.

\subsection{The character quantifier and the precise positivity theorem}

Write \(C_{\mathbb Q}=\mathbb A_{\mathbb Q}^{\times}/\mathbb Q^\times\),
and give its compact norm-one subgroup Haar measure one. Let
\(\mathcal D(C_{\mathbb Q})\) denote finite sums of locally constant
functions on the compact factor times smooth functions of compact
support on \(\mathbb R_+^\times\). Convolution and adjoint use
multiplicative Haar measure:
\[
 (f*g)(x)=\int f(y)g(y^{-1}x)\,d^\times y,\qquad
 f^*(x)=\overline{f(x^{-1})}.
\]
Put
\[
 \widehat f(\chi,z)=
 \int_{\widehat{\mathbb Z}^{\times}}\int_0^\infty
 f(u,r)\overline{\chi(u)}r^z\,\frac{dr}{r}\,du .
\]
Here \(\chi\) labels a primitive Dirichlet character with
Euler coefficients \(\chi(p)\). Its associated Hecke character
has compact-unit part \(\overline\chi\) in the rational
representative convention used here. It is not counted again
each time it is induced to a larger modulus.

\begin{proposition}[AC-T01: the exact rational character range]
\label{AC-T01}
There is a topological group isomorphism
\[
 C_{\mathbb Q}\simeq\widehat{\mathbb Z}^{\times}\times
 \mathbb R_+^\times .
\]
Every continuous \emph{unitary} character of \(C_{\mathbb Q}\) is
\(\chi(u)r^{it}\), for a finite-order Dirichlet character \(\chi\)
and a real number \(t\). Consequently A implies RH for every
unitarily normalized Gr\"ossencharacter of \(\mathbb Q\).
A nonunitary norm twist changes the center of the critical strip;
it must first be removed before real part \(1/2\) is asserted.
\end{proposition}

\begin{proof}
For an idele \(a=(a_\infty,(a_p))\), choose the positive rational
number \(b=\prod_p p^{-v_p(a_p)}\). Then each \(ba_p\) is a unit.
Multiplication by the rational number \(-1\), if necessary, makes
the real component positive. This gives a representative
\((r,u)\) with \(r>0\) and \(u\in\widehat{\mathbb Z}^{\times}\).
A rational number preserving all finite valuations is \(\pm1\);
preserving positivity at infinity then makes it \(1\). The
representative is unique. The construction and its inverse are
continuous, and multiplication of representatives proves the
group assertion. The idele norm is \(r\).

A continuous character of the compact profinite group
\(\widehat{\mathbb Z}^{\times}\) has open kernel. To see this
directly, take an arc about \(1\) in the circle which contains no
nontrivial subgroup, and an open subgroup whose image lies in that
arc. Its image is trivial. Thus the character factors through
\((\mathbb Z/q\mathbb Z)^\times\) for some \(q\).
A continuous unitary character of \(\mathbb R_+^\times\), after
the logarithm, is a character \(x\mapsto e^{itx}\) of
\((\mathbb R,+)\). These observations give the classification.
The corresponding \(L\)-function is \(L(s+it,\overline\chi)\)
in these compact-unit coordinates.
Translation by \(it\) preserves real parts of zeros, so A gives
the assertion. An imprimitive Euler presentation is not a new
Hecke character; the assertion concerns the primitive completed
\(L\)-function, not extra zeros of omitted Euler factors.
\end{proof}

For each primitive \(\chi\), write its nontrivial zeros as
\(\rho=\frac12+i\gamma\), and their multiplicities as \(m_{\chi,\gamma}\).
Under A, all \(\gamma\) are real. Define the centered Weil
functional by its spectral side,
\[
 \mathcal W(f)=
 \sum_{\chi}\sum_\gamma m_{\chi,\gamma}
       \widehat f(\chi,i\gamma).                         \tag{T.1}
\]
For a test in \(\mathcal D\) the outer sum is finite and the inner
sums are absolutely convergent. This is also the centered
geometric distribution given by C99, Appendix II, Theorem 6.
The same identity on the strong Schwartz space used in CCM07 is
CCM07, Theorem 4.16 (label \texttt{varthetagTrH1}) together with
Theorem 6.1 (label \texttt{TraceformulaThm}), with the factor
\(|\cdot|^{-1/2}\) inserted.

\begin{theorem}[AC-T02: the whole rational Weil form and the CCM trace pairing]
\label{AC-T02}
Assuming A,
\[
 \mathcal W(f*f^*)=
 \sum_{\chi,\gamma}m_{\chi,\gamma}
          |\widehat f(\chi,i\gamma)|^2\geq0.             \tag{T.2}
\]
For the relative cyclic cohomology object \(H^1\) of CCM07,
Definition 4.10 (label \texttt{motH1AKCK}), this is precisely
\[
 \operatorname{Tr}\!
 \left(\underline\vartheta_m
       ((\Delta^{-1/2}f)*(\Delta^{-1/2}f^*))\mid H^1\right)
       \geq0,\qquad \Delta^z f(x)=|x|^z f(x).             \tag{T.3}
\]
Equivalently, the trace pairing satisfies
\(\langle h,h^\sharp\rangle\geq0\), where
\(h^\sharp(x)=|x|^{-1}\overline{h(x^{-1})}\).
These assertions hold on CCM07's strong Schwartz domain as well
as on \(\mathcal D\). In the converse direction, positivity of
(T.2) on all compactly supported tests implies A.
\end{theorem}

\begin{proof}
Multiplicative Fourier transformation takes convolution to
multiplication. On a unitary character \((u,r)\mapsto
\overline{\chi(u)}r^{i\gamma}\), it takes \(f^*\) to the conjugate of the
transform of \(f\). Under A every point occurring in (T.1) is
unitary. Substitution gives (T.2) term by term. Absolute
convergence justifies the substitution; equivalently it follows
from the trace-class assertion in the cited source. The same
argument applies to the strong Schwartz space, whose defining
decay makes the spectral trace converge.

The source normalization is
\(\underline\vartheta_m(x)=|x|^{1/2}W(x)\);
see CCM07 (label \texttt{Wvarthetag}) in the proof of
Theorem 4.16. Hence
\(W(f)=\underline\vartheta_m(\Delta^{-1/2}f)\).
Multiplication by \(|x|^{-1/2}\) is an automorphism of convolution,
so this equality converts (T.2) to (T.3).
Setting \(h=\Delta^{-1/2}f\) gives
\(\Delta^{-1/2}f^*=h^\sharp\), proving the second form.
The converse is Weil's positivity criterion in precisely the
form of CCM07, Proposition 6.2
(label \texttt{RHpositTr}) and Corollary 6.3
(label \texttt{RHpositTrsharp}), followed by AC-T01.
Thus no additional unproved positivity assertion is used here.
\end{proof}

The cohomological/nuclear trace identity of CCM07, Theorem 6.1
was already an unconditional theorem; A proves its positive
Hermitian specialization, not its existence. Similarly, the
radical statement and the ability to modify representatives
in CCM07, Lemma 4.17 and Proposition 6.4
(labels \texttt{vanishV}, \texttt{pairTrvanishV}) are
unconditional. A permits quotienting the now positive form
by its full nullspace and completing it. It does not assert
that the previously identified subspace \(\mathcal V\) is
the entire nullspace.

\subsection{What the original spectral realization becomes}

\begin{proposition}[AC-T03: exhaustion, multiplicities, and Jordan blocks]
\label{AC-T03}
For \(k=\mathbb Q\), every character sector in C99,
Theorem III.1, now has as its spectral set \emph{all} the
nontrivial zeros of the corresponding Dirichlet \(L\)-function,
after subtracting \(1/2\). For its Sobolev parameter
\(\delta>1\), a zero of multiplicity \(m\) occurs with the
multiplicity
\[
 m_\delta=\min(m,N_\delta),\qquad
 N_\delta=\max\{n\in\mathbb Z_{\geq1}:n<(1+\delta)/2\}.    \tag{T.4}
\]
Every prescribed finite multiset of zeros is realized with
its full multiplicities by choosing \(\delta\) sufficiently
large. A makes no change to formula (T.4). The stated source
results supply no uniform multiplicity bound independent
of height, and do not establish that this particular
generator becomes skew-adjoint.
\end{proposition}

\begin{proof}
Theorem III.1 states that the spectrum consists exactly of
the critical zeros, with (T.4). A says that every nontrivial
zero is critical, so the spectral-set conclusion follows.
For a finite multiset, let \(M\) be its largest multiplicity
and choose \(\delta>2M-1\). Then \(N_\delta\geq M\).
The supplied theorem contains no assertion that its
multiplicity formula changes when RH holds. Moreover,
the paragraph immediately after C99, III.(30), explicitly
states that a multiple zero produces a nontrivial Jordan
form when \(\delta\) is sufficiently large. A skew-adjoint
operator is normal and has no nontrivial finite-dimensional
Jordan block: if \((D-i\gamma)v=w\ne0\) and
\((D-i\gamma)w=0\), then skew-adjointness gives
\(\|w\|^2=\langle (D-i\gamma)v,w\rangle
=-\langle v,(D-i\gamma)w\rangle=0\), a contradiction.
A is a zero-location assertion, not a simplicity assertion.
Thus the stated skew-adjointness inference would require
further information, not provided by A or by Theorem III.1.
\end{proof}

\begin{theorem}[AC-T04: an explicit positive realization, with its exact scope]
\label{AC-T04}
A permits the following constructions.

First, the pre-Hilbert space
\[
 \mathcal D(C_{\mathbb Q})/\mathcal N,\qquad
 \mathcal N=\{f:\widehat f(\chi,i\gamma)=0
                    \text{ for every }(\chi,\gamma)\},
 \quad
 (f,g)_{\mathcal W}=\mathcal W(f*g^*)
                                                               \tag{T.5}
\]
has a unitary action of \(C_{\mathbb Q}\) by translation.
Its completion is the positive trace-pairing realization.

Second, a multiplicity-preserving trace model is
\[
 \mathcal H_{\rm mult}
  =\bigoplus_{\chi,\gamma}\mathbb C^{m_{\chi,\gamma}},
 \qquad
 H(v_{\chi,\gamma})=(\gamma v_{\chi,\gamma}),\quad
 \operatorname{Dom}H
  =\left\{v:\sum_{\chi,\gamma}
             \gamma^2\|v_{\chi,\gamma}\|^2<\infty\right\}.        \tag{T.6}
\]
The diagonal action
\[
 U(u,r)|_{\mathbb C^{m_{\chi,\gamma}}}
       =\overline{\chi(u)}r^{i\gamma}\operatorname{id}
\]
is unitary and \(H\) is self-adjoint. For finite-character
smooth compactly supported \(f\), the operator \(U(f)\)
is trace class and
\[
 \operatorname{Tr}U(f)=\mathcal W(f).                            \tag{T.7}
\]
The construction (T.6) uses the zeros as input. It is not
an identification with the original adelic cutoff spaces,
nor does (T.5) automatically retain a multiplicity
\(m_{\chi,\gamma}\)-dimensional eigenspace.
\end{theorem}

\begin{proof}
Equation (T.2), polarization, and Cauchy--Schwarz for a
nonnegative sesquilinear form show that quotienting by
its nullspace gives an inner product. The sum-of-squares
formula identifies that nullspace exactly as in (T.5).
Translation multiplies each transform by a unit-modulus
number, so it preserves the inner product and extends
unitarily to the completion. Strong continuity follows
first for a fixed \(f\) by dominated convergence in the
absolutely convergent square sum and then by density.

For (T.6), the direct sum is a Hilbert space. Multiplication
by a real sequence on its maximal square-summability
domain is self-adjoint: for example the diagonal inverses
\((H\pm i)^{-1}\) have coefficients \((\gamma\pm i)^{-1}\),
are bounded, and have ranges in the stated domain.
The action is unitary; dominated convergence in the
square-summable coordinates proves strong continuity.
Integrated against a test \(f\), it is diagonal with
eigenvalue \(\widehat f(\chi,i\gamma)\), repeated
\(m_{\chi,\gamma}\) times. There are only finitely many
character sectors, and rapid vertical decay together
with the usual \(O_\chi(T\log(T+2))\) zero count makes
\(\sum m_{\chi,\gamma}|\widehat f(\chi,i\gamma)|\) finite.
Equivalently, this is the absolute spectral convergence
in C99, Corollary III.2 and CCM07, Theorem 4.16.
This proves trace class and (T.7).

Finally, (T.5) samples a test at a distinct spectral point
only once, with weight \(m_{\chi,\gamma}\). A weight in
a scalar spectral measure does not produce \(m\)
independent vectors at that point. This explains why
(T.5) and the multiplicity-preserving model (T.6) have
different purposes when a zero is multiple.
\end{proof}

\subsection{A fully explicit theorem on every finite window}

Fix a primitive Dirichlet character \(\chi\) of conductor \(q\),
let \(\epsilon\in\{0,1\}\) satisfy \(\chi(-1)=(-1)^\epsilon\),
and write \(\psi=\Gamma'/\Gamma\).
For \(\lambda>1\) set \(L=\log\lambda\), \(I=[-L,L]\).
For \(v\in C_c^\infty((-L,L))\) put
\[
 g(e^t)=v(t),\qquad
 f=g*g^*,\qquad
 F_+v(\tau)=\int_I v(t)e^{i\tau t}\,dt,\qquad
 \ell_\pm(v)=\int_I v(t)e^{\pm t/2}\,dt .
\]
We use the \(L^2\) inner product linear in the second variable.
Let \(T_a\) be compressed translation on \(L^2(I)\):
\[
 (T_a v)(t)=\widetilde v(t-a)\quad(t\in I),
\]
where \(\widetilde v\) is extended by zero outside \(I\).
Thus \(T_a^*=T_{-a}\), and \(T_a=0\) for \(|a|\geq2L\)
(up to the immaterial endpoints). Define the bounded
operator given by a finite sum of translations
\[
 P_{\chi,\lambda}=
 \sum_{\substack{p\ {\rm prime},\,k\geq1\\p^k\leq\lambda^2}}
 \frac{\log p}{p^{k/2}}
 \left(\chi(p)^k T_{-k\log p}
       +\overline{\chi(p)^k}T_{k\log p}\right).                 \tag{T.8}
\]
Each translation on \(L^2(I)\) need not have finite rank.
As usual, \(\chi(p)=0\) at ramified primes.
Let \(\delta_\chi=1\) for the trivial primitive character
(\(q=1\)) and \(0\) otherwise, and set
\[
 a_\chi(\tau)=
 \log(q/\pi)+
 \operatorname{Re}\psi\!\left(\frac{1/2+\epsilon+i\tau}{2}\right).
                                                               \tag{T.9}
\]

\begin{theorem}[AC-T05: the exact semilocal window identity]
\label{AC-T05}
Assuming A, for every \(\lambda>1\) and every such \(v\),
\[
\begin{split}
 q_{\chi,\lambda}(v)
 &:=
 \int_{\mathbb R}a_\chi(\tau)|F_+v(\tau)|^2\,
                    \frac{d\tau}{2\pi}
 -\langle v,P_{\chi,\lambda}v\rangle\\
 &\hspace{6mm}
 +2\delta_\chi\operatorname{Re}
            \bigl(\ell_+(v)\overline{\ell_-(v)}\bigr)\\
 &=\sum_{\rho=1/2+i\gamma}
               m_{\chi,\gamma}|F_+v(\gamma)|^2\ \geq0 .
\end{split}\tag{T.10}
\]
Only the prime powers \(p^k\leq\lambda^2\) occur.
For \(\chi=1\), imposing either \(\ell_+(v)=0\) or
\(\ell_-(v)=0\) removes the two polar terms. For a nontrivial
primitive character there are no polar terms to remove.
Positivity of all these forms, for all \(\lambda\) and all
primitive \(\chi\), is equivalent to A.
\end{theorem}

\begin{proof}
The completed Dirichlet function has archimedean factor
\[
 (q/\pi)^{(s+\epsilon)/2}
       \Gamma((s+\epsilon)/2).
\]
Twice the real part of its logarithmic derivative on
\(s=1/2+i\tau\) is (T.9). Applying the explicit formula
C99, Appendix II, Theorem 6 in the \(\chi\)-sector
therefore gives, for an arbitrary smooth compactly
supported \(f\) in the centered normalization,
\[
\begin{split}
 \sum_\rho m_\rho\int_0^\infty
             f(x)x^{\rho-1/2}\,\frac{dx}{x}
={}&
 \delta_\chi\left(\int f(x)x^{1/2}\frac{dx}{x}
             +\int f(x)x^{-1/2}\frac{dx}{x}\right)\\
 &+\int_{\mathbb R}a_\chi(\tau)
           \left(\int f(x)x^{i\tau}\frac{dx}{x}\right)
                   \frac{d\tau}{2\pi}\\
 &-\sum_{p,k\geq1}\frac{\log p}{p^{k/2}}
           \bigl(\chi(p)^k f(p^k)
             +\overline{\chi(p)^k}f(p^{-k})\bigr).
\end{split}\tag{T.11}
\]
For clarity, the conductor term comes from
\(2\,d\log(q^{s/2})/ds=\log q\); the real gamma term comes
from pairing the two functional-equation contours; and
the prime terms follow from
\(-L'/L(s,\chi)=\sum_{p,k}(\log p)\chi(p)^k p^{-ks}\)
on \(\operatorname{Re}s>1\).
These are the local components in the cited general
explicit formula. The \(\chi=1\) version, with the same
signs, is also printed in CC20, Appendix B, equations
\texttt{bombieriexplicit}--\texttt{burnolexplicit1}.
The gamma factor has no zeros; its trivial-zero
cancellations are already incorporated in (T.11).

Now \(f(e^a)=\int v(t)\overline{v(t-a)}dt\), so
\(\operatorname{supp}f\subset[\lambda^{-2},\lambda^2]\),
and
\[
 f(e^a)=\langle v,T_{-a}v\rangle,\qquad
 \int f(x)x^{i\tau}\frac{dx}{x}=|F_+v(\tau)|^2 .
\]
Consequently the prime sum is exactly (T.8); terms outside
the stated range vanish. The two polar transforms sum
to \(2\operatorname{Re}(\ell_+\overline{\ell_-})\).
Under A, \(\rho-1/2=i\gamma\) with \(\gamma\) real.
Substitution in (T.11) proves the identity and inequality.

For the converse, every compact logarithmic support is
contained in some interval \([-L,L]\). The hypotheses
therefore give the centered Weil positivity in each
character sector. The sectors are orthogonal under
convolution and finite sums give all tests in
\(\mathcal D(C_{\mathbb Q})\). AC-T02 and Weil's
criterion imply A. The endpoint constraints in the
trivial sector can also be imposed without losing the
criterion: CC20, Proposition 46
(label \texttt{mainprop}), proves exactly this
finite-codimension version.
\end{proof}

\paragraph{Prior result and the extension made here.}
For the trivial character, the closed Weil form and
compact-resolvent operator are already established in
Connes--Consani--Moscovici, \emph{Zeta spectral triples}
(arXiv:2511.22755), \S3, Proposition 3.3
(label \texttt{Hilbert}), Proposition 3.5
(label \texttt{compres}), and Theorem 3.6
(label \texttt{thmsmallest}). The following proof
records the required form for every primitive
Dirichlet character, including odd and complex
characters, and the consequence of A.

\begin{theorem}[AC-T06: closed Dirichlet window operators]
\label{AC-T06}
The forms (T.10) have a canonical closed realization on
\(L^2(I)\). Under A their associated self-adjoint
operators \(A_{\chi,\lambda}\) are nonnegative and have
compact resolvent. In particular their spectra consist
of nonnegative eigenvalues of finite multiplicity tending
to \(+\infty\). Their kernels are precisely
\[
 \ker A_{\chi,\lambda}
 =\{v\in\mathcal D_{\chi,\lambda}:
                 F_+v(\gamma)=0
          \text{ for all zeros }1/2+i\gamma
          \text{ of }L(s,\chi)\},                             \tag{T.12}
\]
where the form domain \(\mathcal D_{\chi,\lambda}\) is
specified in the proof. Every finite matrix obtained by
polarizing (T.10) on test vectors is positive semidefinite.
The same statements hold for the form restricted to
either closed polar-annihilating hyperplane
\(\ker\ell_+\) or \(\ker\ell_-\).
\end{theorem}

\begin{proof}
The digamma function has no poles on the vertical lines
in (T.9). Stirling's formula gives
\[
 a_\chi(\tau)=\log|\tau|+\log(q/(2\pi))+O_\chi(|\tau|^{-2})
       \quad(|\tau|\longrightarrow\infty).
\]
It is continuous, bounded below, and tends to infinity.
Choose \(c_\chi\) such that \(b_\chi=a_\chi+c_\chi\geq1\).
This weight is comparable above and below to
\(1+\log(2+|\tau|)\). Define
\[
 \mathcal D_{\chi,\lambda}
 =\left\{v\in L^2(I):
  \int b_\chi(\tau)|F_+v(\tau)|^2\,\frac{d\tau}{2\pi}
                      <\infty\right\}.                       \tag{T.13}
\]
The positive weighted Fourier form is closed: extension
by zero identifies \(L^2(I)\) with a closed subspace of
\(L^2(\mathbb R)\), and multiplication by
\(b_\chi^{1/2}\) is a closed operator in Fourier space.

Smooth tests in the open interval are a form core.
Indeed first replace \(v(t)\) by
\(v_r(t)=r^{-1/2}v(t/r)\) with \(r<1\) tending to \(1\).
Dilation is strongly continuous in the weighted Fourier
norm. This follows by approximation by compactly supported
Fourier functions and the uniform bound on
\(b_\chi(\tau/r)/b_\chi(\tau)\) for \(1/2\leq r\leq1\).
The support of \(v_r\) lies in the interior of \(I\).
Convolving \(v_r\) with a smooth compactly supported
approximate identity of sufficiently small radius
preserves support in \(I\). In Fourier space it multiplies
by a uniformly bounded function tending pointwise to
\(1\); dominated convergence gives form-norm convergence.
This proves the core assertion.

The operator (T.8) is bounded and self-adjoint, with
\(\|P_{\chi,\lambda}\|\leq
2\sum_{p^k\leq\lambda^2}(\log p)p^{-k/2}\).
The polar term is a bounded Hermitian form because
\(\ell_\pm\) are bounded functionals on \(L^2(I)\).
Hence adding these bounded terms to the closed weighted
Fourier form gives a closed semibounded form on (T.13).
On the core it is (T.10), so A makes it nonnegative
on its whole domain. The representation theorem for
closed nonnegative forms defines
\(A_{\chi,\lambda}\) uniquely.

Here is the compactness argument. A form-norm bounded
set has uniformly small Fourier tails:
\[
 \int_{|\tau|>R}|F_+v(\tau)|^2\frac{d\tau}{2\pi}
 \leq
 \frac{C}{\inf_{|\tau|>R}b_\chi(\tau)}
 \longrightarrow0.
\]
The Fourier transform from \(L^2(I)\) to \(L^2([-R,R])\)
has square-integrable kernel \(e^{i\tau t}\), and is
therefore compact. Low-frequency compactness and the
uniform tail bound show that the embedding of the form
domain into \(L^2(I)\) is compact. The form resolvent
maps bounded sets of \(L^2(I)\) to form-norm bounded
sets, proving compact resolvent.

To extend the sum-of-squares formula, approximate
\(v\in\mathcal D_{\chi,\lambda}\) by core vectors \(v_n\).
The sequences
\((\sqrt{m_{\chi,\gamma}}F_+v_n(\gamma))_\gamma\)
are Cauchy in \(\ell^2\), because their squared distance
is \(q_{\chi,\lambda}(v_n-v_m)\). Each coordinate
evaluation is continuous on \(L^2(I)\), by
Cauchy--Schwarz. The limit is therefore the corresponding
sequence for \(v\), and (T.10) extends to the form domain.
A vector belongs to the kernel of a nonnegative operator
if and only if its form value is zero. This gives (T.12).

Restriction to either hyperplane is again a closed form.
Core density in that hyperplane follows by taking smooth
approximants and subtracting their small moment error
times one fixed smooth test whose corresponding moment
is nonzero. The compactness argument persists under
restriction. Finally a polarized Gram matrix is
positive semidefinite because its quadratic form is
\(q_{\chi,\lambda}(\sum_j c_jv_j)\geq0\).
\end{proof}

These constructed window operators have the Weil form
as their quadratic form. Their spectral conclusion, strengthened to strict positivity
in AC-T13 below, does not identify their eigenvalues with
the zero ordinates. No assertion
of a uniform spectral gap as \(\lambda\to\infty\), nor
a convergence rate for any different spectral-triple
approximation, is included in AC-T06.

\subsection{What the semilocal compressions satisfy}

Let \(\vartheta(r)\xi(x)=r^{-1/2}\xi(r^{-1}x)\) on
\(L^2(\mathbb R)_{\rm ev}\), and let \(\mathbf S\) be
the Sonin projection onto the even functions which,
together with their additive Fourier transforms,
vanish on \([-1,1]\). CC20, Theorem 27
(label \texttt{devil}), is the unconditional identity
\[
 \operatorname{Tr}(\vartheta(f)\mathbf S)
      =W_\infty(f)+E(f),\qquad
 E(f)=\int_0^\infty f(r^{-1})\varepsilon(r)\frac{dr}{r}.          \tag{T.14}
\]
Here \(W_\infty=-W_{\mathbb R}\) in that paper's unitary
normalization, and its explicit convergent prolate
expression is
\[
 \varepsilon(r)=
 \sum_n\frac{\lambda(n)}{\sqrt{1-\lambda(n)^2}}
       \langle \xi_n,\vartheta(r^{-1})\zeta_n\rangle
 \quad(r\geq1),\qquad \varepsilon(r^{-1})=\varepsilon(r).
                                                               \tag{T.15}
\]
All vectors and eigenvalues in this formula are defined
in CC20, \S4, Proposition 25 (label \texttt{devil0});
the convergence treatment is Appendix F, including
Lemma 49 (label \texttt{lemesti}).

\begin{theorem}[AC-T07: corrected Sonin and semilocal trace inequalities]
\label{AC-T07}
Let \(g\in C_c^\infty([\lambda^{-1},\lambda])\),
\(f=g*g^*\), and assume either
\(\int g(r)r^{1/2}d^\times r=0\) or
\(\int g(r)r^{-1/2}d^\times r=0\).
Assuming A,
\[
\begin{split}
 0\leq \sum_{\rho=1/2+i\gamma}m_\rho
                  |\widehat g_+(\gamma)|^2
 ={}&\operatorname{Tr}
          (\vartheta(g)\mathbf S\vartheta(g)^*)-E(f)\\
 &-\sum_{p^k\leq\lambda^2}
      \frac{\log p}{p^{k/2}}
             \bigl(f(p^k)+f(p^{-k})\bigr),
                                                               \\
 \widehat g_+(\tau)&=\int g(r)r^{i\tau}d^\times r .
\end{split}\tag{T.16}
\]
Thus the Sonin trace dominates the explicit sum of
the prolate remainder and these finitely many prime
contributions, for every window.

More generally, let \(S\) contain infinity, the
ramified primes of a character sector, and the primes
occurring in its test. Put \(v(t)=g(e^t)\) and
\(f_\chi(u,r)=\chi(u)f(r)\), viewed on \(C_S\)
after quotienting by the outside-\(S\) compact units.
In the unitary normalization of
C99, Theorem VII.4, the finite part of the semilocal
cutoff trace in that sector is
\[
 \operatorname{FP}_{\Lambda\to\infty}
 \operatorname{Tr}(R_\Lambda\vartheta_S(f_\chi))
     =B_\chi(f)-q_{\chi,\lambda}(v),                           \tag{T.17}
\]
where the subtracted divergent term is
\(2f(1)\log'\Lambda\), and \(B_\chi(f)\) is the polar
term in (T.10). It is therefore at most \(B_\chi(f)\);
on the polar-annihilating subspace it is nonpositive.
This concerns the finite part, not the sign of each
finite-\(\Lambda\) cutoff trace.
\end{theorem}

\begin{proof}
For the trivial character, (T.11), with the polar
terms removed, gives
\(q=W_\infty(f)-\sum_p W_p(f)\), where
\[
 W_p(f)=(\log p)\sum_{k\geq1}p^{-k/2}
                  (f(p^k)+f(p^{-k})).
\]
Insert (T.14). For \(f=g*g^*\), the trace in (T.14)
equals \(\operatorname{Tr}
(\vartheta(g)\mathbf S\vartheta(g)^*)\) by the
trace-class and cyclicity assertions proved there.
This proves the equality in (T.16); A supplies
nonnegativity through (T.10).

For (T.17), write the original nonunitary scaling
action as \(U(a)=|a|^{1/2}\vartheta_S(a)\).
In C99, Theorem VII.4, substitute the test
\(h(a)=|a|^{-1/2}f_\chi(a)\), after averaging the
irrelevant compact-unit coordinates. Then
\(U(h)=\vartheta_S(f_\chi)\) and \(h(1)=f(1)\).
The finite part is exactly the sum of local
geometric terms. Since the chosen \(S\) contains
all nonzero terms for this test, the explicit
formula (T.11) identifies this local sum with
\(B_\chi(f)-q_{\chi,\lambda}(v)\). This proves
(T.17) and its inequalities. The theorem uses
\(R_\Lambda=\widehat P_\Lambda P_\Lambda\), which
need not be an orthogonal projection, so the
argument makes no pointwise-in-\(\Lambda\)
positivity claim.
\end{proof}

CC20's original small-window theorem is stronger than
the bare inequality \(W_\infty(g*g^*)\geq0\) in a
different direction. Theorem 1 assumes support in
\([2^{-1/2},2^{1/2}]\) and two transform vanishings.
Theorem 44 (label \texttt{mainthmfine}) proves
\[
 W_\infty(g*g^*)\geq
 \operatorname{Tr}(\vartheta(g)\mathbf S\vartheta(g)^*)
                   -c|\widehat g(0)|^2,\qquad 13<c<17,
                                                               \tag{T.18}
\]
under its specified one endpoint vanishing. These
were already source theorems. Their quantitative
constant depends on the prolate estimates and
numerical ingredients in CC20; this run does not
certify those numerical computations anew.
A proves (T.16) on all windows, but does not make
the explicit \(E\) and prime terms zero. The quantitative bound in (T.18) still uses the
cited prolate estimates. CC20, Remark 45
(label \texttt{remlowc}), gives the source
obstruction to removing the positive constant
under those hypotheses.

\begin{proposition}[AC-T08: compression positivity does not give a restricted flow]
\label{AC-T08}
For any unitary representation \(U\) and orthogonal
projection \(P\), the compression
\(\Phi(f)=PU(f)P|_{P\mathcal H}\) satisfies
\[
 \Phi(g*g^*)-\Phi(g)\Phi(g)^*
   =PU(g)(1-P)U(g)^*P\geq0.                                  \tag{T.19}
\]
The kernel \((x,y)\mapsto PU(x^{-1}y)P\) is positive
definite. These facts are unconditional. The
compressed family \(PU(x)P\) need not be a
representation, so positivity alone does not
produce the spectral theorem for a restricted
scaling flow.
\end{proposition}

\begin{proof}
Use \(U(g*g^*)=U(g)U(g)^*\), insert \(1=P+(1-P)\),
and subtract. For the kernel, vectors
\(\xi_j\in P\mathcal H\) give
\[
 \sum_{i,j}\langle\xi_i,PU(x_i^{-1}x_j)P\xi_j\rangle
  =\left\|\sum_j U(x_j)\xi_j\right\|^2\geq0.
\]
For a counterexample to the representation assertion,
take \(U(t)=\operatorname{diag}(e^{it},e^{-it})\)
and project onto \((1,1)/\sqrt2\). The compression
is the scalar \(\cos t\), and
\(\cos(t+s)\ne\cos t\cos s\) in general.
This verifies the failure even in a finite-dimensional
setting with no domain issues. For a compression of
an unbounded generator there are, in addition,
domain and self-adjointness questions; (T.19)
does not answer them.
\end{proof}

\begin{proposition}[AC-T09: a compactly supported counterexample to global archimedean positivity]
\label{AC-T09}
The claim that the archimedean term \(W_\infty\)
alone is nonnegative on every window, even after
both polar moments vanish, is false. This failure
is compatible with A.
\end{proposition}

\begin{proof}
In logarithmic coordinates the archimedean form is
the Fourier multiplier
\[
 a_1(\tau)=\operatorname{Re}\psi(1/4+i\tau/2)-\log\pi.
\]
The reflection and duplication formulas for
\(\Gamma\) give
\[
 a_1(0)=-\gamma_{\rm E}-\frac\pi2-3\log2-\log\pi<0.
\]
For example this follows from
\(\psi(1/2)=-\gamma_{\rm E}-2\log2\),
\(\psi(3/4)-\psi(1/4)=\pi\), and
\(\psi(1/4)+\psi(3/4)=2\psi(1/2)-2\log2\).

Choose a nonzero real \(\phi\in C_c^\infty((-1,1))\),
put \(\phi_T(t)=T^{-1/2}\phi(t/T)\), and
\[
 v_T=(\partial_t^2-1/4)\phi_T,\qquad g_T(e^t)=v_T(t).
\]
Integration by parts proves
\(\ell_+(v_T)=\ell_-(v_T)=0\).
The support of \(g_T\) lies in \([e^{-T},e^T]\).
Its archimedean form is
\[
 \int_{\mathbb R} a_1(\eta/T)
    ((\eta/T)^2+1/4)^2
       |F_+\phi(\eta)|^2\frac{d\eta}{2\pi}.
\]
As \(T\to\infty\) this tends to
\(a_1(0)\|\phi\|_2^2/16<0\).
Dominated convergence is valid since \(a_1(\tau)\)
has logarithmic growth and \(F_+\phi\) decays
faster than any power. Thus a sufficiently large
finite \(T\) gives a smooth compactly supported
counterexample. The full form in (T.10) contains
prime terms as well. Consequently there is no
contradiction with A.
\end{proof}

\subsection{The 1999 global cutoff: source scope and a conductor-level reduction}

Theorem VIII.5 of C99 explicitly assumes positive
characteristic. Its implication from the cutoff
trace to positivity uses the exact inclusion
\(E(B_{\Lambda,0})\subset S_\Lambda\), C99 VIII.(23),
and its reverse implication uses Lemma VIII.3.
At the end of \S VIII, equations (30)--(33), C99
discusses the number-field case, restricts to
\(k=\mathbb Q\), uses even real prolate functions,
and states that this gives the analogues of
Lemma 1, Theorem 5, and Lemma 3 for the Riemann
zeta function. Thus, \emph{if that stated
number-field analogue is taken as a cited
theorem}, A gives its zeta cutoff formula:
\[
 \operatorname{Tr}(Q_\Lambda U(h))
 =2h(1)\log'\Lambda+
   \sum_v\int_{\mathbb Q_v^\times}'
        \frac{h(u^{-1})}{|1-u|}\,d^\times u+o(1).               \tag{T.20}
\]
This is a consequence of the asserted analogue,
with the prolate cutoff meant there. It is not a
derivation of a full all-character number-field
cutoff theorem from the positive-characteristic
Theorem VIII.5. The introduction's all-character
formulation and the nuclear formula of CCM07
do not supply the missing cutoff estimate.

\begin{proposition}[AC-T10: why the literal function-field cutoff cannot be copied]
\label{AC-T10}
A nonzero function on \(\mathbb R\) cannot have
both itself and its additive Fourier transform
compactly supported. Hence the number-field
extension must use an approximate/prolate
replacement for the literal simultaneous-support
cutoff in Theorem VIII.5.
\end{proposition}

\begin{proof}
If \(f\in L^2(\mathbb R)\) has compact support,
then it is in \(L^1\), and
\[
 \widehat f(z)=\int f(x)e^{-2\pi i xz}\,dx
\]
extends to an entire function by differentiation
under the compactly supported integral.
If its restriction to the real axis also has
compact support, continuity makes it vanish
on a real open interval. The identity theorem
makes it identically zero, and injectivity
of Fourier transformation gives \(f=0\).
This obstruction is explicitly noted in
C99, \S VIII just before equation (31).
\end{proof}

\begin{proposition}[AC-T11: finite conductor, Gauss phase, and parity]
\label{AC-T11}
Let \(\chi\) be a primitive nontrivial Dirichlet
character of conductor \(q\), and extend it
by zero to nonunits modulo \(q\). On the finite
adeles put
\[
 \varphi_\chi(x)=
   \sum_{a\bmod q}\chi(a)\,
                     1_{a+q\widehat{\mathbb Z}}(x).
                                                               \tag{T.21}
\]
Choose the standard finite additive character
so that for \(y=b/q+\widehat{\mathbb Z}\) its
value on the residue \(a\) is \(e^{2\pi iab/q}\).
With \(\operatorname{vol}(\widehat{\mathbb Z})=1\),
\[
 \widehat\varphi_\chi(y)=
 \begin{cases}
 q^{-1}\tau(\chi)\overline{\chi(b)},&
           y=b/q+\widehat{\mathbb Z},\\
 0,&y\notin q^{-1}\widehat{\mathbb Z},
 \end{cases}
 \qquad
 \tau(\chi)=\sum_{a\bmod q}\chi(a)e^{2\pi ia/q}.                 \tag{T.22}
\]
In particular it vanishes unless \(b\) is a unit
modulo \(q\), and \(|\tau(\chi)|=\sqrt q\).
For an archimedean vector of parity
\(\chi(-1)=(-1)^\epsilon\), simultaneous
\emph{nominal} adelic norm cutoffs \(\Lambda\)
are therefore satisfied by the uniform archimedean rectangle
\[
 |x|\leq\Lambda,\qquad
 |\xi|\leq\Lambda/q.                                         \tag{T.23}
\]
Unitary dilation transforms this to the symmetric
rectangle of side parameter \(\Lambda/\sqrt q\).
Pointwise on the chosen finite tensor factor, this rectangle
is a sufficient cutoff: fibers with smaller finite norm permit
a larger real coordinate.  AC-T47 subsequently passes through
the actual semilocal quotient and proves that the minimal
conductor set has exactly this raw pair and metric.
AC-T48 compares the genuine angle cutoffs for each additional
fixed finite set of places.  Thus the pointwise observation
is not a remaining obstruction to that rational source
dictionary.
Even and odd prolate sectors handle
\(\epsilon=0,1\), respectively. The scalar
Gauss phase disappears from the geometry of the
two orthogonal projections.
\end{proposition}

\begin{proof}
The Fourier transform of \(1_{a+q\widehat{\mathbb Z}}\)
is \(q^{-1}\) times the indicated additive
character on \(q^{-1}\widehat{\mathbb Z}\),
and zero outside that annihilator. This gives
the finite sum in (T.22).
For a unit \(b\), changing variables \(a\mapsto
b^{-1}a\) turns the sum into
\(\overline{\chi(b)}\tau(\chi)\).
For nonunit \(b\), choose a prime divisor \(p\)
of \((b,q)\). Primitivity implies that the
character is nontrivial on the kernel of the
reduction from the units modulo \(q\) to the
units modulo \(q/p\). Choose \(u\) in that
kernel with \(\chi(u)\ne1\). The additive
character in the sum is unchanged by
\(a\mapsto ua\), since \(p\mid b\);
the multiplicative character is multiplied by
\(\chi(u)\). Thus the sum is zero.
Finite Fourier Plancherel now gives
\[
 q\sum_{a\bmod q}|\chi(a)|^2
   =\sum_{b\bmod q}
        |\tau(\chi)\overline{\chi(b)}|^2,
\]
so \(q\varphi(q)=|\tau(\chi)|^2\varphi(q)\).

The support of (T.21) has finite norm at most
\(1\), while the support of (T.22) has finite
norm at most \(q\); at every prime dividing
\(q\) the latter valuation is exactly the
negative conductor exponent. Combining with
the real factor gives (T.23). If
\(D_r h(x)=r^{-1/2}h(x/r)\), then \(D_r\)
multiplies the time support by \(r\) and
divides the Fourier support by \(r\). Take
\(r=\sqrt q\) and side
\(\Lambda/\sqrt q\).
Fourier transformation preserves parity.
Multiplication of a Fourier transform by a
unit scalar has no effect on its support
projection. These facts prove the last
assertions.
\end{proof}

There is an explicit finite-dimensional version
of the remaining analytic problem. For a finite
span of compactly supported prolate vectors
\(\psi_1,\ldots,\psi_N\) of the required parity
in (T.23), define
\[
 e_j(x)=x^{1/2}\sum_{n\geq1}\chi(n)\psi_j(nx)
       \quad(x>0),\quad
 G_{ij}=\langle e_i,e_j\rangle_{L^2(dx/x)},\quad
 M(h)_{ij}=\langle e_i,V(h)e_j\rangle .                         \tag{T.24}
\]
For fixed nontrivial \(\chi\), its periodic partial
sums are bounded. Summation by parts therefore
bounds the sum defining \(e_j(x)/\sqrt x\)
uniformly as \(x\downarrow0\), for the bounded-
variation prolate cutoff vectors. It vanishes
for \(x>\Lambda\). Thus \(e_j\in L^2(dx/x)\).
Their Mellin transforms in a right half-plane
are \(L(s,\chi)\) times the Mellin transforms
of the \(\psi_j\); uniqueness proves linear
independence. Hence \(G\) is positive definite.
The orthogonal projection \(Q_{\chi,\Lambda}\)
onto their span is concretely
\(E G^{-1}E^*\), where \(Ec=\sum_jc_je_j\), and
\[
 \operatorname{Tr}(Q_{\chi,\Lambda}V(h))
           =\operatorname{tr}(G^{-1}M(h)).                    \tag{T.25}
\]
The expected prolate threshold is the one in
C99, VIII.(32), rescaled as in (T.23):
indices of the required parity up to
\(4\Lambda^2/q\). For the trivial character
one must also account for the two endpoint
conditions and their polar correction.

\paragraph{The original cutoff target and its subsequent resolution.}
For the prolate image spaces just described
(and a cutoff convention whose threshold errors
are controlled), the target was
\[
 2h(1)\log\Lambda
 -\operatorname{tr}(G_{\chi,\Lambda}^{-1}M_{\chi,\Lambda}(h))
 \longrightarrow
 \sum_\gamma m_{\chi,\gamma}\widehat h_+(\gamma)               \tag{T.26}
\]
for every smooth compactly supported \(h\),
with the polar correction in the trivial
sector. Before A, the proposed limit is the
sum of the harmonic measures in C99,
Lemma VIII.3; A collapses these measures to
Dirac masses if this limit theorem has been
established. The finite Fourier calculation
AC-T11 settles the conductor and parity
part. It does not estimate the inverse Gram
matrices or the tails of the \(E_\chi\)-images.

Two attempted shortcuts fail at identifiable
steps. First, exact simultaneous support
fails by AC-T10. Second, a prolate vector
has a nonzero Fourier tail, so its \(E_\chi\)
image need not vanish below \(\Lambda^{-1}\).
The exact inclusion \(Q'_{\Lambda,0}\leq
S_\Lambda\) in C99, VIII.(23) is therefore
not available without an error estimate.
Orthogonality of prolate vectors for the
real \(L^2\) norm does not give orthogonality
for the \(E_\chi\) norm: the matrix \(G\)
is essential. Expanding its double arithmetic
sum and interchanging integrals term by
term requires convergence not supplied by
bounded character partial sums. Thus
(T.26) was not proved by these initial calculations.
AC-T28--AC-T41 subsequently prove it for every
fixed nonprincipal primitive Dirichlet sector at
the full original threshold.  AC-T42--AC-T44
prove the principal two-constraint version, and
AC-T45 supplies its explicit full polar completion.
The supplied positive-characteristic theorem alone
does not supply those analytic arguments.
AC-T47 supplies the exact rational source quotient metric
and the minimal-set time-band pair, and AC-T48 proves
tested-trace invariance for every fixed semilocal angle
cutoff.  A simultaneous growing-place-set limit or a
different unspecified quantization is outside those
statements.

\subsection{Geometric language and the statements that remain}

For a finite-order \(\chi\), the arithmetic
object is a rank-one finite-monodromy
\(\ell\)-adic local system on
\(\operatorname{Spec}\mathbb Z[1/q\ell]\),
or its Artin motive with coefficients
containing the character values. Its
unramified Frobenius eigenvalues already
have complex absolute value \(1\).
A supplies an analytic statement about the
completed \(L\)-function; it does not newly
prove that elementary weight-zero fact.
Equations (T.2) and (T.10) are positive
pairings of test functions attached to
these \(L\)-functions.

CCM07, Definition 4.10 constructs a cokernel
in a category of cyclic modules and
equation \texttt{TorH1mot} applies a derived
functor to it. This is the precise
categorical location of the source's
cohomological terminology. The supplied
construction does not identify it with
finite-dimensional \(\ell\)-adic
cohomology of a proper scheme, or with
a specified pure motive. AC-T04 adds
a positive Hilbert completion and an
explicit spectral trace model. Neither
construction supplies an algebraic
square, a \(K\)-theory comparison map,
or algebraic correspondences realizing
the pairing.

For the intersection notation of CCM07,
Proposition 7.2 (label \texttt{PosTrRHint}),
A gives the stated numerical inequality
in the real test-function convention.
For complex tests, the endpoint term
written there as \(2\widehat f(0)
\widehat f(1)\) must be replaced by
\[
 2\operatorname{Re}
      \bigl(\widehat f(0)\overline{\widehat f(1)}\bigr).
\]
This follows directly by transforming
\(f*f^\sharp\), as in (T.10).
It is a Hermitian correction, not a
new geometric intersection construction.

Beyond the now proved rational source tensor
formula, further questions concern uniformity
for a simultaneous growing-place-set limit;
an identification of its prolate image
spaces with the positive trace-pairing
completion; and any proposed semilocal
Sonin domination stronger than (T.16),
which requires a separate estimate for
the explicit remainder. The local
trace formula C99, Theorem V.3, the
semilocal trace formula Theorem VII.4,
the CCM07 nuclear trace formula, and
CC20's already proved small-window
theorems are not newly proved by A.

\subsection{The inherited slow-band trace theorem and a sharper window tool}

\begin{theorem}[AC-T12: A closes the MAS\_ARTIN slow-band arithmetic gate over \(\mathbb Q\)]
\label{AC-T12}
Fix \(0<\alpha<1/2\) and \(s\geq0\). Use the
actual source operators of MAS\_ARTIN checkpoint 10,
\texttt{modules/full\_test\_slow\_band.tex},
MA-FT-001 and MA-FT-002, in
\[
 H_{\alpha,s}
 =L^2\!\left(\mathbb R,
         e^{2\alpha|x|}(1+x^2)^{s/2}\,dx\right).
\]
In the rational trivial sector and every primitive
nonprincipal Dirichlet sector, A gives, for every
\(g\in C_c^\infty(\mathbb R)\) with \(g(0)=0\),
\[
 L_{\alpha,s}(g)=\sum_\rho m_\rho
        \int_{\mathbb R}g(t)e^{i\gamma_\rho t}\,dt,
 \qquad \Delta_{\alpha,s}(g)=0.                               \tag{T.27}
\]
In particular the paired actual local and global
annular-cutoff limits of MA-FT.(7) agree whenever
\(X,Y\to\infty\) with \(Y/X\to0\), and with the
fixed common annular profile specified there.
The source-thin version additionally retains
its condition \(\delta X\to0\).
This is true at every positive strip width,
with full multiplicities and for every \(s\geq0\).
\end{theorem}

\begin{proof}
The source defines \(d=1/2-\alpha\), \(u=1/2+\alpha\)
and retained multiplicities \(r_\rho^+\) and
\(r_\rho^{\rm st}\), with possible deficiencies
\(\delta_\rho^\pm\), in MA-FT.(2).
Under A every zero has real part \(1/2\), strictly
between \(d\) and \(u\). Hence
\(r_\rho^+=r_\rho^{\rm st}=m_\rho\) and
\(\delta_\rho^-=\delta_\rho^+=0\); no zero lies
on a strip boundary, so the boundary jet
parameter \(b_s\) has no effect.
Substitute these values in MA-FT.(5)--(7).
All deficiency sums disappear, giving (T.27)
and the stated actual operator-limit
identities. The support, profile, and
path hypotheses remain precisely those
of MA-FT-001, rather than a new order
of limits.
\end{proof}

The proof of MA-FT-001/002 was read in full
for this use. Its ambient trace-class,
near-norm-one, and local-quotient results
are retained as the explicitly named
background dependencies of that proof;
they have not all been re-proved in
this section. The statement that A
annihilates its two deficiency divisors
is checked directly above.

This supplies another actual trace realization.
It leaves a concrete comparison problem before
one obtains C99's prolate cutoff: the source
uses positive \(\alpha\), one-sided
Paley--Wiener theta-image closures and
a slow-band joint limit; C99 uses the
unweighted simultaneous time/frequency
prolate geometry. A comparison must
control the trace after the Hilbert
space identification, including passage
\(\alpha\downarrow0\), and must restore
the identity contact term. Indeed
\(g(0)=0\) excludes every nonzero test
of the form \(g=v*v^*\), since
\((v*v^*)(0)=\|v\|_2^2\).
Equality on this test class alone cannot
determine a distribution's coefficient
of \(\delta_0\). Two distributions
differing by \(c\delta_0\) agree on
every such \(g\). This is the precise
missing normalization, not a claim
that the source theorem is wrong.

\begin{corollary}[AC-T13: strictly positive spectrum in every fixed Dirichlet window]
\label{AC-T13}
Assuming A, \(\ker A_{\chi,\lambda}=0\) for
every primitive Dirichlet character and every
finite \(\lambda>1\). Consequently
\(\min\operatorname{Spec}A_{\chi,\lambda}>0\).
This is a statement for each fixed window;
it supplies no uniform lower bound as the
window grows.
\end{corollary}

\begin{proof}
Under GRH, Carneiro--Chandee--Milinovich,
\emph{A note on the zeros of zeta and
\(L\)-functions}, arXiv:1503.00955v1,
Theorem 5, (4.4), and Example 8, give
\(S(t,\chi)=O_\chi(\log|t|/\log\log|t|)\).
The jump at a zero bounds its multiplicity
by the same order. The
zero-count identity (4.1) in that source gives
\(\asymp_\chi T\log T\) zeros with multiplicity; hence there
are \(\gg_\chi T\log\log T\) distinct
ordinates with \(|\gamma|\leq T\).

If \(v\) belonged to (T.12), its entire
transform \(F_+v(z)\) would vanish at all
these ordinates. Cauchy--Schwarz gives
\(|F_+v(z)|\leq\sqrt{2L}\|v\|_2e^{L|\operatorname{Im}z|}\).
For a nonzero such entire function,
Jensen's formula about any point at
which it does not vanish bounds its
number of zeros in \(|z|\leq T\) by
\(O_v(T)\): apply Jensen at radius
\(2T+O(1)\), and use the displayed
growth bound. This contradicts the
distinct-ordinate count. Thus \(v=0\).
Compact resolvent and nonnegativity
in AC-T06 then imply that the least
eigenvalue is strictly positive.
\end{proof}

For the trivial character this is the
same strictness argument as AC-S03
in the spectral-triple part of this
work; AC-T13 supplies its primitive
Dirichlet extension.

\begin{proposition}[AC-T14: the exact norm of a compressed prime translation]
\label{AC-T14}
On \(L^2([-L,L])\), for \(a>0\) and \(z\in\mathbb C\),
\[
 \|zT_a+\overline zT_{-a}\|=
 \begin{cases}
  2|z|\cos\!\left(\dfrac{\pi}
             {\lceil 2L/a\rceil+1}\right),&0<a<2L,\\
  0,&a\geq2L.
 \end{cases}                                                 \tag{T.28}
\]
In particular if \(L\leq a<2L\), the norm
is \(|z|\), rather than the triangle
bound \(2|z|\). For any finite sum
of prime translations, summing the
individual quantities (T.28) gives
a valid operator-norm bound.
\end{proposition}

\begin{proof}
Write almost every point of the interval
uniquely as
\(t=-L+r+ja\), with \(0\leq r<a\) and
\(j=0,\ldots,N(r)-1\), omitting values
outside the interval. This is a unitary
direct-integral decomposition into
finite coordinate chains: integration
in \(t\) becomes integration in \(r\)
followed by summation in \(j\).
On each chain \(T_a\) is the one-step
shift with zero boundary condition.
If \(0<a<2L\), the essential maximum
chain length is
\(N_{\max}=\lceil2L/a\rceil\).
When \(2L/a\) is an integer, the
additional endpoint occurs only
for a measure-zero set of \(r\);
there is no extra coordinate in
the essential maximum.

A diagonal unitary on a chain removes
the phase of \(z\). The remaining
matrix is \(|z|\) times the adjacency
matrix of a path with \(N(r)\)
vertices. Its eigenvectors have
coordinates
\(\sin(jk\pi/(N(r)+1))\),
\(j=1,\ldots,N(r)\), and eigenvalues
\(2|z|\cos(k\pi/(N(r)+1))\),
\(k=1,\ldots,N(r)\), as direct
substitution in the two-term
recurrence verifies. Its norm is
therefore
\(2|z|\cos(\pi/(N(r)+1))\).
Taking the essential supremum in
the direct integral gives (T.28).
For \(a\geq2L\) both compressed
translations vanish almost
everywhere. The last assertion
is the operator-norm triangle
inequality.
\end{proof}
